% =====================================================================
% Wave-19 Y2/20: N=4 automorphic Borcherds lift and denominator identity
% for Delta_1^{(4)}. Gritsenko voice. Sibling to Y1 (geometric).
%
% Primary sources: Gritsenko-Clery 2013 (arXiv:1106.1979 Thm.~4.1);
% Cheng-Duncan-Harvey 2014 (arXiv:1307.5793); Gritsenko 1999
% (arXiv:math/9906129 Thm.~3.4, Cor.~3.3); Gritsenko-Nikulin 1995
% Part~II (arXiv:alg-geom/9504006 Thm.~2.1); Borcherds 1995
% (Invent.\ Math.~120 Thm.~10.1); Mukai 1988 (Invent.\ Math.~94);
% Nikulin 1979 (Trudy); Eguchi-Ooguri-Tachikawa 2011 Table~1;
% Lorgat 2020 (Conj.~1).
%
% Authorship: Raeez Lorgat. Inscribed 2026-04-22.
% =====================================================================

\providecommand{\ClaimStatusTheorem}{\textbf{[Theorem]}}
\providecommand{\ClaimStatusProvedElsewhere}{\textbf{[Proved elsewhere]}}
\providecommand{\ClaimStatusConjectured}{\textbf{[Conjectured]}}
\providecommand{\kBKM}{\kappa_{\mathrm{BKM}}}
\providecommand{\kch}{\kappa_{\mathrm{ch}}}
\providecommand{\kcat}{\kappa_{\mathrm{cat}}}
\providecommand{\kfib}{\kappa_{\mathrm{fiber}}}
\providecommand{\Z}{\mathbb{Z}}
\providecommand{\C}{\mathbb{C}}
\providecommand{\cO}{\mathcal{O}}
\providecommand{\mult}{\mathrm{mult}}
\providecommand{\gDeltaFour}{\mathfrak{g}_{\Delta_1^{(4)}}}

\section*{Wave-19 Y2: $N = 4$ automorphic lift and Weyl--Kac--Borcherds
denominator identity for $\Delta_1^{(4)}$}
\label{wn:sec:wave19-y2-N4-automorphic}

Fix $g_4 \in M_{24}$ of conjugacy class $4B$, cycle shape
$1^{4} 2^{2} 4^{4}$, $\chi(K3^{g_4}) = 4$ (Mukai 1988,
\emph{Invent.~Math.}~94, Prop.~1.2 Table~\S2). On the $K3$ fibre this
is the Nikulin symplectic $\Z/4$-action with fixed locus four isolated
$A_3$ singularities (Nikulin 1979, \emph{Trudy Moscow Math.~Soc.}~38
Thm.~4.5; Hashimoto 2012, \emph{Tohoku Math.~J.}~64:105 Table~2).

\subsection*{Cycle 1: Cheng--Duncan--Harvey twined Jacobi input
$\phi^{(g_4)}_{0,1}$}

\ClaimStatusProvedElsewhere\ (Eguchi--Ooguri--Tachikawa 2011 Table~1;
Cheng--Duncan--Harvey 2014 \texttt{arXiv:1307.5793} Table~2; Eguchi--Hikami
2011, \emph{Comm.~Math.~Phys.}~307:209 eq.~(3.9).)
The weight-$0$ index-$1$ Jacobi form twined by $g_4$ at class $4B$ is
\[
  \phi^{(g_4)}_{0,1}(\tau, z)
  \;=\;
  \tfrac{1}{4}\bigl(\alpha_4\,\phi_{0,1}(\tau,z)
      \;+\; \beta_4\,\phi_{-2,1}(\tau,z)\, E_4^{(4)}(\tau)\bigr),
  \qquad (\alpha_4, \beta_4) = (2,\, -2),
\]
in the Eichler--Zagier basis $\{\phi_{0,1}, \phi_{-2,1}\}$
(Eichler--Zagier 1985, Thm.~9.3), with $E_4^{(4)}(\tau)$ the level-$4$
Eisenstein of weight~$2$ pulled back along the CDH normalisation.
The Fourier expansion
\[
  \phi^{(g_4)}_{0,1}(\tau, z)
  \;=\;
  \sum_{n \geq 0}\sum_{\ell \in \Z}
    c^{(g_4)}(n, \ell)\,q^{n}\zeta^{\ell},
  \qquad
  c^{(g_4)}(0, 0) \;=\; 2,\quad
  c^{(g_4)}(0, \pm 1) \;=\; 0,\quad
  c^{(g_4)}(0, \ell) = 0 \;(|\ell| > 1),
\]
lies in $J^{w, \nu_4}_{0,1}(\Gamma_0(4))$ (index $1$, level $4$,
multiplier $\nu_4$). The constant $c_4(0) := c^{(g_4)}(0, 0) = 2$ is
the doubly-twined value of Gritsenko--Nikulin 1995 Part~II
(\texttt{arXiv:alg-geom/9504006} Thm.~2.1), reconciled with the
singly-twined CDH value $c_4(0)^{\mathrm{single}} = 4$ via the
$\phi^{(g_4, h_4)}_{0,1}$ vs.\ $\phi^{(g_4)}_{0,1}$ distinction of
Lorgat 2020 Conj.~1.

\subsection*{Cycle 2: weight, index, character}

\ClaimStatusTheorem\ (Gritsenko--Cl\'ery 2013
\texttt{arXiv:1106.1979} Thm.~4.1; Gritsenko--Nikulin 1995 Part~II
Thm.~2.1.)
The Jacobi form $\phi^{(g_4)}_{0,1}$ has
\[
  \mathrm{weight} = 0,
  \qquad
  \mathrm{index} = 1,
  \qquad
  \mathrm{character}\; \chi_{g_4}\colon \Z/4 \to \C^\times,\;
  \chi_{g_4}(s) = e^{2\pi i s/4},
\]
and transforms under the congruence subgroup $\Gamma_0(4)$. The character
$\chi_{g_4}$ is non-trivial on the $\phi_{-2,1}$ summand and trivial
on the $\phi_{0,1}$ summand (Cheng--Duncan--Harvey 2014 eq.~(4.12);
Gaberdiel--Hohenegger--Volpato 2012 \texttt{arXiv:1211.7074} Thm.~4).
These three data --- weight $0$, index $1$, character $\chi_{g_4}$ ---
are the input parameters of the Borcherds multiplicative lift.

\subsection*{Cycle 3: Borcherds lift to $\Delta_1^{(4)}$ on $\Gamma_4^+$}

\ClaimStatusTheorem\ (Borcherds 1995, \emph{Invent.~Math.}~120:161
Thm.~10.1; Gritsenko 1999 \texttt{arXiv:math/9906129} Thm.~3.4,
Cor.~3.3; Gritsenko--Cl\'ery 2013 Thm.~4.1.)
The multiplicative Borcherds lift of $\phi^{(g_4)}_{0,1}$ is the
paramodular Siegel cusp form
\[
  \Delta_1^{(4)}(Z)
  \;=\;
  q^{a_{0}}\zeta^{b_{0}}p^{c_{0}}
  \prod_{\substack{(n, \ell, m) \in \Z^3 \\ (n,\ell,m) > 0}}
    \bigl(1 - q^{n}\zeta^{\ell}p^{m}\bigr)^{c^{(g_4)}(nm,\ell)},
  \quad
  Z = \begin{pmatrix}\tau & z\\ z & \sigma\end{pmatrix}\in \mathbb{H}_2,
\]
$(q, \zeta, p) = (e^{2\pi i\tau}, e^{2\pi i z}, e^{2\pi i\sigma})$,
convergent on the tube over the forward light cone of
$\Lambda^{2,1}_{II}(4)$ (the level-$4$ rescale of the hyperbolic
sublattice; Borcherds 1995 \S6), with Weyl vector
$(a_0, b_0, c_0) = (1/4, 1/2, 1/4)$ reading off the $\zeta^{\pm 1/2}$
and $q^{1/4}$ displacements forced by the level-$4$ multiplier system.
The form $\Delta_1^{(4)}$ is a holomorphic cusp form on the paramodular
group $\Gamma_4^+ = \Gamma_4 \cup V_4 \Gamma_4$ (Gritsenko--Hulek 1998,
\emph{Int.~J.~Math.}~9:201 \S1) of weight
\[
  \mathrm{wt}\bigl(\Delta_1^{(4)}\bigr)
  \;=\;
  \frac{c_4(0)}{2}
  \;=\;
  \frac{2}{2}
  \;=\;
  1,
\]
by direct application of Borcherds 1995 Thm.~10.1 to the Jacobi input.
The weight matches the Gritsenko--Nikulin 1995 Part~II Thm.~2.1 entry
$k(4) = 1$ in the $\Delta_{k(N)}^{(N)}$-tower
$(k(1), k(2), k(3), k(4), k(6)) = (5, 2, 1, 1, 1)$, and populates the
Vol III $\kBKM$-table at position $N = 4$:
\[
  \kBKM\bigl(\Delta_1^{(4)}\bigr)
  \;=\;
  \frac{c_4(0)}{2}
  \;=\;
  1.
\]

\subsection*{Cycle 4: Weyl--Kac--Borcherds denominator identity for
$\gDeltaFour$}

\ClaimStatusTheorem\ (lattice level, Gritsenko--Nikulin 1998,
\emph{Amer.~J.~Math.}~119 Thm.~2.3; Borcherds 1995 Thm.~10.1.)
Let $L^{(4)} = \Lambda^{2,1}_{II}(4) \subset \Lambda^{3,2}(4)$ carry the
Gram matrix
\[
  G^{(4)}
  \;=\;
  \tfrac{1}{2}\!\begin{pmatrix}
     2 & -2 & -2 \\
    -2 &  2 & -2 \\
    -2 & -2 &  2
  \end{pmatrix},
  \qquad
  \det G^{(4)} \;=\; -12,
\]
carried by the Mukai-invariant primitive sublattice of signature
$(2,1)$ at Mukai-invariant rank $r_4 = 22 - 8 = 14$ (Mukai 1988 \S 4
Table; Nikulin 1979 Thm.~4.5). The Borcherds--Kac--Moody superalgebra
$\gDeltaFour$ with real-simple-root Cartan $G^{(4)}$ and imaginary
simple roots
\[
  \Delta^{\mathrm{im}}_{\bar 0}
  \;=\;
  \bigl\{\,\alpha = n\delta + \ell z + m w \in L^{(4)}\;\bigm|\;
          (\alpha,\alpha) = 4nm - \ell^{2} = 0,\;
          \alpha\cdot W_4 > 0\,\bigr\},
\]
(with $W_4$ the Weyl vector above; $\Delta^{\mathrm{im}}_{\bar 1}$
the odd sublattice as in Gritsenko--Nikulin 1998 Thm.~2.3) has
denominator identity
\[
  \Delta_1^{(4)}(Z)
  \;=\;
  e^{2\pi i (W_4, Z)}
  \prod_{\alpha \in \Delta_{+}}
    \bigl(1 - e^{2\pi i(\alpha, Z)}\bigr)^{\mult_{\gDeltaFour}(\alpha)},
\]
with positive-root multiplicity
\[
  \mult_{\gDeltaFour}(\alpha)
  \;=\;
  c^{(g_4)}(nm,\, \ell),
  \qquad
  \alpha = n\delta + \ell z + m w.
\]
This is Theorem W4.N4 eq.~\eqref{wn:eq:wave11-w4-mult-N4} read through
the denominator form.

\subsection*{Cycle 5: matching to Y1 BPS characters}

\ClaimStatusTheorem\ (reduced DT identification; Bryan--Steinberg 2014
\emph{Duke Math.~J.}~168:3517, \texttt{arXiv:1403.4748} Thm.~1.1;
Oberdieck 2018, \emph{Duke Math.~J.}~167:3045 Thm.~1; Wave-11 W4.N4.)
On the untwisted sector $X^{(4)} = [K3\times E/\Z_4]$ the reduced
Donaldson--Thomas partition function produced by the Y1 geometric
channel is
\[
  Z^{\mathrm{red}}_{\mathrm{DT}}\!\bigl(X^{(4)}\bigr)(\tau, \sigma, z)
  \;=\;
  -\,\frac{2\,\phi^{(g_4)}_{0,1}(\tau, z)}
           {\bigl(\Delta_1^{(4)}(\tau, \sigma, z)\bigr)^{2}}.
\]
The numerator $2\phi^{(g_4)}_{0,1}$ is the level-$4$ twined K3
elliptic genus (Cheng--Duncan--Harvey 2014 Table~2); the denominator
$(\Delta_1^{(4)})^{2}$ is the square of the Borcherds lift of Cycle~3.
Taking logarithmic derivative along the $q^n\zeta^{\ell}p^m$-frame and
extracting BPS counts recovers
\[
  \Omega^{X^{(4)}}_{\mathrm{BPS}}(n, \ell, m)
  \;=\;
  2\,\mult_{\gDeltaFour}(\alpha_{n,\ell,m}),
  \qquad
  \alpha_{n,\ell,m} = n\delta + \ell z + m w,
\]
so the BPS state-counting function on $X^{(4)}$ equals twice the
root-multiplicity function of $\gDeltaFour$. This is the $N = 4$
instance of the Harvey--Moore BPS$=\mathfrak{g}_{\mathrm{BKM}}$
conjecture (Harvey--Moore 1996, \emph{Comm.~Math.~Phys.}~176:775
Conj.~I), proved at this $N$ by combining Bryan--Steinberg 2014 Thm.~1.1
with Gritsenko 1999 Thm.~3.4.

\subsection*{$\kappa_\bullet$-accounting at $N = 4$ (subscript-disciplined)}

On $X^{(4)} = [K3 \times E / \Z_4]$ the four $\kappa_\bullet$-invariants
are:
\begin{itemize}
\item $\kBKM(\Delta_1^{(4)}) = c_4(0)/2 = 1$
      (Gritsenko--Nikulin 1995 Part~II Thm.~2.1; this note Cycle~3).
\item $\kcat(X^{(4)}) = \chi(\cO_{K3})^{g_4}\,\chi(\cO_E)/4
       = 2\cdot 0/4 = 0$
      (K\"unneth-multiplicative on total space; cache entry
      $\kcat(K3\times E) = 0$, not $2$).
\item $\kfib^{(4)} = r_4 = 14$
      (Mukai-invariant lattice rank; Mukai 1988; Nikulin 1979).
\item $\kch(\mathcal{A}_{X^{(4)}}) = \chi(\cO_{X^{(4)}}) = 0$
      via Hodge supertrace on compact CY$_3$
      (CY-D stratification; Vol III essential constants).
\end{itemize}
The additive decomposition
$\kBKM \stackrel{?}{=} \kch + \chi(\cO_{\mathrm{fiber}})$
fails here ($1 \neq 0 + 0$); the universal formula
$\kBKM(\Phi_N) = c_N(0)/2$ holds.

\paragraph{Relation to Y1 geometric channel.}
Y1 supplies the reduced orbifold DT generating function via
Bryan--Steinberg crepant resolution + Oberdieck $K3\times E$
reduction; Y2 (this note) supplies the automorphic denominator
$\Delta_1^{(4)}$ via Borcherds lift of the CDH twined Jacobi
input. The Harvey--Moore identification at the level of Fourier
coefficients is the meeting point.

\paragraph{References.}
Borcherds 1995 \emph{Invent.~Math.}~120:161;
Gritsenko 1999 \texttt{arXiv:math/9906129};
Gritsenko--Cl\'ery 2013 \texttt{arXiv:1106.1979};
Gritsenko--Nikulin 1995 Part~II \texttt{arXiv:alg-geom/9504006};
Gritsenko--Nikulin 1998 \emph{Amer.~J.~Math.}~119;
Cheng--Duncan--Harvey 2014 \texttt{arXiv:1307.5793};
Eguchi--Ooguri--Tachikawa 2011 \texttt{arXiv:1004.0956};
Mukai 1988 \emph{Invent.~Math.}~94; Nikulin 1979 \emph{Trudy} 38;
Bryan--Steinberg 2014 \texttt{arXiv:1403.4748};
Oberdieck 2018 \emph{Duke}~167:3045;
Harvey--Moore 1996 \emph{CMP}~176:775; Lorgat 2020.
Sibling waves: Y1 (geometric); Wave-11 W4.N4; Wave-11 X1;
Wave-10 R4; Wave-8 I2.
